\section{审判监督程序（再审程序）}
\begin{enumerate}
  \item 4 个问题。
        \begin{enumerate}
          \item 再审分为几个独立的程序，分别解决什么问题。2 个（逻辑出发点）。\yi 启动程序。\er 重新审理程序。
          \item 当事人申请再审\zx{不停止原判决的执行}；法院裁定再审时应当裁定中止原判的执行。
                例外。\al{\xhx{医疗费}、赡养费、扶养费、抚育费、抚恤金、\xhx{劳动报酬}} \yjt \zd{继续执行}。
                请问这两句话冲突吗，它们的关系是什么。选项看到裁定中止执行的时候提高警惕性！！
          \item 法院在裁定再审时，能不能一并裁定撤销原判。作出新的判决才能撤销原判。
          \item 在整个再审程序当中，原生效判决的效力如何。
        \end{enumerate}
  \item 时间线。% https://q.uiver.app/?q=WzAsOCxbMCwyLCJPXFxcXCJdLFsyLDIsIkFcXFxc6KOB5a6a5YaN5a6hIl0sWzQsMiwiQlxcXFzmlrDliKTlhrMiXSxbMiwwLCLkuK3mraLmiafooYwiXSxbNCwwLCLmkqTplIDljp/liKQiXSxbMCwwXSxbMSwxLCLlkK/liqgiXSxbMywxLCLlrqHnkIYiXSxbMCwxLCIiLDAseyJzdHlsZSI6eyJoZWFkIjp7Im5hbWUiOiJub25lIn19fV0sWzEsMiwiIiwwLHsic3R5bGUiOnsiaGVhZCI6eyJuYW1lIjoibm9uZSJ9fX1dLFsxLDMsIiIsMCx7InN0eWxlIjp7ImJvZHkiOnsibmFtZSI6ImRhc2hlZCJ9LCJoZWFkIjp7Im5hbWUiOiJub25lIn19fV0sWzIsNCwiIiwwLHsic3R5bGUiOnsiYm9keSI6eyJuYW1lIjoiZGFzaGVkIn0sImhlYWQiOnsibmFtZSI6Im5vbmUifX19XSxbMCw1LCIiLDIseyJzdHlsZSI6eyJib2R5Ijp7Im5hbWUiOiJkYXNoZWQifSwiaGVhZCI6eyJuYW1lIjoibm9uZSJ9fX1dXQ==
\[
  \begin{tikzcd}[row sep=tiny]
    {} && \txt{中止执行} && \txt{撤销原判} \\
    & \txt{启动} && \txt{审理} \\
    O && \txt{A\\裁定再审} && \txt{B\\新判决}
    \arrow[no head, from=3-1, to=3-3]
    \arrow[no head, from=3-3, to=3-5]
    \arrow[dashed, no head, from=3-3, to=1-3]
    \arrow[dashed, no head, from=3-5, to=1-5]
    \arrow[dashed, no head, from=3-1, to=1-1]
  \end{tikzcd}
\]

  \item 对象。
        \begin{enumerate}
          \item 已生效且确有错误的判决、裁定（不予受理、驳回起诉）、调解书。
          \item 解除婚姻关系不准再审。\al{\xhx{离婚诉讼}不允许再审 \cuo{离婚诉讼 \bdy 解除婚姻关系。财产分割、子女抚养}}。
        \end{enumerate}
  \item 启动。
        \begin{enumerate}
          \item 法院。本院和上级法院。
          \item 检察院。
                \begin{enumerate}
                  \item （再审）抗诉。\yi 上抗下，抗诉书同级提。\er 再审检察建议同级提，报上级检察院备案。
                  \item 抗诉必再审。接受抗诉的法院应当在收到抗诉书之日起 30 日内直接裁定再审
                        （不能审查，说你错了你就错了，没有说不的权利）。
                \end{enumerate}
          \item 当事人。要有\zd{再审利益/理由}（13 个，\bs{民诉法 207}）。
                \begin{enumerate}
                  \item 自知道 \hz 应当知道裁判文书生效之日起的 6 个月内提出。
                  \item \zd{原则找上级法院。当事人一方人数众多 \hz 当事人双方都是公民可以找原审法院}（方便当事人）。\al{甲起诉乙人身损害}。
                  \item 一方找上级，一方找原审 \yjt 无法协商 \yjt 原审受理 \yjt 3 月审查。
                  \item \zx{只有 2 个结果。裁定再审 \hz 裁定驳回再审申请}。\al{不予受理 \cuo}。
                  \item 只有 2 次机会，不能同时提。法院 \yjt 不处理 \hz 对再审判决不满意 \yjt 检察院申诉。断后性、终局性。
                        \begin{enumerate}
                          \item 驳回再审申请。
                          \item 逾期未对再审申请作出裁定。
                          \item 再审判决、裁定仍有明显错误。
                        \end{enumerate}
                \end{enumerate}
          \item 案外人。\yi 第三人 \yjt 三撤。\er 本应作为原审共同原、被告 \yjt 申请再审。
        \end{enumerate}
  \item 重新审理。
        \begin{enumerate}
          \item 法院。谁启动谁来审，谁惹麻烦谁收场。
          \item 检察院。抗诉的由接受抗诉的法院提审 \hz 指令下一级法院（事实、证据问题抗诉，再审审理过的除外 \yjt 屡教不改）。
                \al{广州中院一审 \yjt 自己启动再审 \yjt 检抗诉 \yjt 高院 \yjt | 广州中院}。
          \item 当事人申请的再审。
                \begin{enumerate}
                  \item \bs{向谁申请谁来审。不能低于原审法院级别}。
                        \zd{因当事人申请而裁定再审的案件由中级以上法院审理}。
                  \item 例外。\bs{一审是基层法院作出，向中院/基层法院申请 \yjt 只能中院审（提审）/基层法院审}！！
                        不区分是当事人还是检的申请。
                \end{enumerate}
          \item 必须合议庭。
          \item 程序。\zx{原生效裁判一审作出按一审，二审按二审，凡是提审统统按二审}。做题先看是不是提审。
          \item \bs{再审范围有限原则}！！\zx{增加、变更诉讼请求，提出反诉 \yjt 不予审理}（\al{判决、组织调解 \cuo}）。例外。
                \begin{enumerate}
                  \item 人。
                        \begin{enumerate}
                          \item 漏人。原审被缺席判决 \yjt 剥夺了当事人的诉讼权利（辩论权）。
                          \item 加人。追加了新的当事人 \yjt 平等原则。追加了 = 只看结果，不管谁加进来的。
                        \end{enumerate}
                  \item 事。无法通过另诉解决 \yjt 山穷水尽。
                  \item 物。标的灭失 \yjt 原请求无法继续实现 \yjt 审下去无意义。
                \end{enumerate}
          \item 看到增加请求先看是几审！！
          \item \zx{4 种结案方式}。
                \begin{enumerate}
                  \item 裁判。跟原判一样的维持。不一样的，撤销 + 改判。三撤只有是你的才会改。
                  \item 调解、撤回起诉同二审。
                  \item 终结再审。人死无继承、当事人和解并履行完毕、未经授权申请再审、已在程序重复申请。
                        \begin{enumerate}
                          \item 启动阶段 \yjt 裁定终结再审审查。
                          \item 重新审理阶段 \yjt 裁定终结再审程序。
                        \end{enumerate}
                \end{enumerate}
          \item 再审发回原审。原审 = 原一审、原二审。注意区别选项指代啥。
        \end{enumerate}
        \begin{kk}{法院对大力公司提出的再审请求如何处理？为什么？}
          \bs{啥程序，怎么判}。
          \begin{enumerate}
            \item 高院\zx{应当裁定再审}，\zx{提审}或者\zx{指令}原审中院或者其他中院审理。
            \item 本案的再审应当\zx{适用二审程序}重新审理。
            \item 对于\zx{遗漏解除合同这一诉讼请求}，再审法院可以进行调解，调解不成，撤销原一审、二审判决，发回一审法院重新审理。
          \end{enumerate}
        \end{kk}
\end{enumerate}
